\section{Original full-quartet critical values and every ordered tensor sum}
\label{sec:quartet-critical-values}

Fix the original offcritical point and the original equal full order
\[
 \rho=\tfrac12+\delta+i\gamma,\qquad
 0<\delta<\tfrac12,\quad \gamma>0,\quad m\in\mathbb Z_{\geq1}.
 \tag{QCV.1}
\]
The order is the full multiplicity in the arithmetic divisor; no
simple-root replacement is made. Keep the ordered roots
\[
 (r_1,r_2,r_3,r_4)=(\rho,\overline\rho,1-\rho,1-\overline\rho),
 \quad h(s)=\prod_{j=1}^4(s-r_j)^m,
 \quad \Phi_h(s)=\int_0^s h(t)\,dt.
 \tag{QCV.2}
\]
The last integral means the unique polynomial antiderivative with
\(\Phi_h(0)=0\). All paths give the same polynomial. In particular,
the constant in this phase is fixed, and is not removed below.

\subsection{The exact polynomial, primitive and local leading coefficients}

The translation \(z=s-\tfrac12\) gives the explicit identity
\[
 \begin{split}
 f(z)&=h(z+\tfrac12)
   =\bigl[((z-\delta)^2+\gamma^2)((z+\delta)^2+\gamma^2)\bigr]^m\\
   &=\bigl[z^4+2(\gamma^2-\delta^2)z^2
                      +(\delta^2+\gamma^2)^2\bigr]^m,
 \qquad F(z)=\int_0^z f(x)\,dx,\qquad A=F(\tfrac12).
 \end{split}
 \tag{QCV.3}
\]
These equations preserve the original \(s\)-coordinate through the
inverse map \(s=z+\tfrac12\). The polynomial \(f\) is even with real
coefficients. It is strictly positive on the real axis: each quadratic
factor in its first expression is positive because \(\gamma>0\).
Thus \(F\) is odd with real coefficients, and \(A>0\). Direct
differentiation and evaluation at zero prove
\[
 \Phi_h(s)=F(s-\tfrac12)+A,
 \quad \Phi_h(1)=2A,
 \quad \Phi_h(1-s)=2A-\Phi_h(s),
 \quad \Phi_h(\overline s)=\overline{\Phi_h(s)}.
 \tag{QCV.4}
\]
In particular the positive constant \(A\) remains in every critical
value and becomes \(kA\) in an ordered \(k\)-fold tensor phase.

For a completely finite expression put
\[
 a_j=\sum_{\substack{r,t\geq0\\r+t\leq m\\r+2t=j}}
 \frac{m!}{r!t!(m-r-t)!}
 [2(\gamma^2-\delta^2)]^r
 (\delta^2+\gamma^2)^{2(m-r-t)}\quad(0\leq j\leq2m).
 \tag{QCV.5}
\]
The multinomial theorem, followed by termwise polynomial integration,
gives
\[
 f(z)=\sum_{j=0}^{2m}a_jz^{2j},\qquad
 F(z)=\sum_{j=0}^{2m}\frac{a_jz^{2j+1}}{2j+1},\qquad
 A=\sum_{j=0}^{2m}\frac{a_j}{(2j+1)2^{2j+1}}.
 \tag{QCV.6}
\]
There is no convergence issue in these finite sums, including when
\(\gamma<\delta\) and some \(a_j\) have negative signs.

Define the real numbers \(U,V\) by the exact evaluation
\[
 U+iV=F(\delta+i\gamma)
     =\sum_{j=0}^{2m}\frac{a_j(\delta+i\gamma)^{2j+1}}{2j+1}.
 \tag{QCV.7}
\]
Oddness and conjugation now give, in the unchanged order (QCV.2),
\[
 (\Phi_h(r_1),\Phi_h(r_2),\Phi_h(r_3),\Phi_h(r_4))
   =(A+U+iV,A+U-iV,A-U-iV,A-U+iV).
 \tag{QCV.8}
\]
The local coefficient at the first root, and all its partners, are
\[
 \begin{gathered}
 b_1=\left.\frac{h(s)}{(s-r_1)^m}\right|_{s=r_1}
       =[8i\delta\gamma(\delta+i\gamma)]^m\ne0,\qquad
 b_2=\overline{b_1},\quad b_3=(-1)^m b_1,
       \quad b_4=(-1)^m\overline{b_1},\\
 \Phi_h(r_j+y)=\Phi_h(r_j)
          +\frac{b_j}{m+1}y^{m+1}+y^{m+2}R_j(y),
       \qquad R_j\in\mathbb C[y].
 \end{gathered}
 \tag{QCV.9}
\]
Indeed the three differences at \(r_1\) are \(2i\gamma\),
\(2\delta+2i\gamma\), \(2\delta\), whose product gives \(b_1\).
Conjugation gives \(b_2\); the substitution \(s\mapsto1-s\) gives
the sign \((-1)^m\) in \(b_3,b_4\). Expanding the polynomial \(h\)
and integrating proves the last line. This is the local coefficient
of the full polynomial (QCV.2), and not of the smaller polynomial
\((s-r_j)^m\).

\subsection{A positive sum of squares proves that the rectangle does not collapse}

For \(d>0\) put \(P_d(x)=(d^2-x^2)^m\). The exact identities are
\[
 \boxed{\begin{aligned}
 2U&=\sum_{j=0}^{m}\frac{(2\gamma)^{2j}}{(2j)!}
         \int_{-\delta}^{\delta}\bigl(P_\delta^{(j)}(x)\bigr)^2\,dx,\\
 2V&=\sum_{j=0}^{m}\frac{(2\delta)^{2j}}{(2j)!}
         \int_{-\gamma}^{\gamma}\bigl(P_\gamma^{(j)}(x)\bigr)^2\,dx.
 \end{aligned}}
 \qquad U>0,\quad V>0.
 \tag{QCV.10}
\]
Here every derivative is the ordinary derivative of the displayed
real polynomial. These are finite identities, with all factorials
and interval lengths retained.

To prove the first identity, integrate \(f\) along the segment from
\(-\delta+i\gamma\) to \(\delta+i\gamma\). Its primitive difference
is \(F(\delta+i\gamma)-F(-\delta+i\gamma)=2U\). On this segment
\[
 \begin{split}
 f(x+i\gamma)
  &=\bigl[(x^2-\delta^2)((x+2i\gamma)^2-\delta^2)\bigr]^m
    =P_\delta(x)P_\delta(x+2i\gamma),\\
 P_\delta(x+2i\gamma)
  &=\sum_{r=0}^{2m}\frac{(2i\gamma)^r}{r!}P_\delta^{(r)}(x).
 \end{split}
 \tag{QCV.11}
\]
An odd derivative of the even polynomial \(P_\delta\) is odd,
so its product with \(P_\delta\) has zero integral on the symmetric
interval. For \(0\leq j\leq m\), integrate by parts \(j\) times:
\[
 \int_{-\delta}^{\delta}P_\delta P_\delta^{(2j)}\,dx
   =(-1)^j\int_{-\delta}^{\delta}(P_\delta^{(j)})^2\,dx.
 \tag{QCV.12}
\]
At the \(r\)-th integration, every boundary term contains a derivative
of \(P_\delta\) of order at most \(j-1<m\). It vanishes at both
endpoints because \(P_\delta\) has a zero of order \(m\) there.
The factor \((2i\gamma)^{2j}=(-1)^j(2\gamma)^{2j}\) cancels
the sign in (QCV.12), proving the first identity in (QCV.10).
Its \(j=0\) integral is strictly positive because \(P_\delta(0)>0\).

Write \(F_{d,e}\) for the polynomial in (QCV.3) with parameters
\((d,e)\) and the same \(m\). Its defining polynomials satisfy
\[
 f_{\delta,\gamma}(iw)=f_{\gamma,\delta}(w),\qquad
 F_{\delta,\gamma}(iw)=iF_{\gamma,\delta}(w).
 \tag{QCV.13}
\]
The first identity follows by substituting \(iw\) in (QCV.3);
the second follows by differentiating and comparing the zero
constant terms. At \(w=\gamma-i\delta\), it gives
\(U_{\delta,\gamma}+iV_{\delta,\gamma}
 =i\overline{U_{\gamma,\delta}+iV_{\gamma,\delta}}\).
Thus \(V_{\delta,\gamma}=U_{\gamma,\delta}>0\), and the first
identity with exchanged parameters proves the second one in
(QCV.10). No restriction on \(\gamma/\delta\) was used.

If desired each integral in (QCV.10) has the fully rational finite
evaluation
\[
 \int_{-d}^{d}(P_d^{(j)}(x))^2\,dx
 =2d^{4m-2j+1}
 \sum_{r,t=\lceil j/2\rceil}^{m}
 \frac{(-1)^{r+t}\binom mr\binom mt(2r)!(2t)!}
 {(2r-j)!(2t-j)![2(r+t-j)+1]}.
 \tag{QCV.14}
\]
Expand \(P_d(x)=\sum_{r=0}^m(-1)^r\binom mr d^{2(m-r)}x^{2r}\),
differentiate \(j\) times, multiply, and integrate the even monomials.
The exponent of \(d\) is \(4m-2j+1\) in every term. Positivity is
proved by the integral in (QCV.10); no sign assertion about the
individual summands of (QCV.14) is needed.

\subsection{An exact real-part bound, retaining the original constant}

The following bound on the same critical values is elementary:
\[
 \begin{gathered}
 0<U<\delta\,[\delta^2(\delta^2+4\gamma^2)]^m,\\
 \gamma\geq\delta\ \Longrightarrow\
 A>\tfrac12(\delta^2+\gamma^2)^{2m},\qquad
 \frac{U}{A}<2\delta
 \left(\frac{\delta^2(\delta^2+4\gamma^2)}
                 {(\delta^2+\gamma^2)^2}\right)^m.
 \end{gathered}
 \tag{QCV.15}
\]
For its proof \(F(i\gamma)\) is purely imaginary, since \(F\) is
odd with real coefficients. Therefore
\[
 U=\operatorname{Re}\int_0^\delta f(x+i\gamma)\,dx.
 \tag{QCV.16}
\]
Put \(t=\delta^2-x^2\), so \(0\leq t\leq\delta^2\). Direct
multiplication gives the exact squared modulus of the expression
inside the \(m\)-th power:
\[
 \begin{split}
 B(x)&=(x^2-\delta^2)
              [x^2-\delta^2-4\gamma^2+4i\gamma x],\\
 |B(x)|^2&=t^2[t^2-8\gamma^2t+16\gamma^4+16\gamma^2\delta^2],\\
 \frac{d}{dt}|B|^2
   &=4t[t^2-6\gamma^2t+8\gamma^4+8\gamma^2\delta^2]>0
          \qquad(0<t\leq\delta^2).
 \end{split}
 \tag{QCV.17}
\]
The strict sign follows from
\(8\gamma^2\delta^2-6\gamma^2t\geq2\gamma^2\delta^2>0\).
Thus \(|B(x)|\leq\delta^2(\delta^2+4\gamma^2)\), with equality
only at \(x=0\). Integrating its \(m\)-th power on a positive-length
interval proves the strict upper bound for \(U\) in (QCV.15),
using the already proved \(U>0\).
When \(\gamma\geq\delta\), the second expression for \(f(x)\)
in (QCV.3) is at least \((\delta^2+\gamma^2)^{2m}\), strictly
larger for \(x>0\). Integrating from \(0\) to \(1/2\) proves the
bound for \(A\), and division proves the ratio bound.

For the original range \(0<\delta<1/2\) and \(\gamma\geq2\), put
\(r=\delta^2/\gamma^2<1/16\). Since
\[
 q(r)=\frac{r(r+4)}{(1+r)^2},\qquad
 q'(r)=\frac{4-2r}{(1+r)^3}>0\quad(0\leq r\leq1/16),
 \qquad q(1/16)=\frac{65}{289},
 \tag{QCV.18}
\]
equation (QCV.15), including its \(2\delta\) factor, implies
\[
 0<U<A\left(\frac{65}{289}\right)^m<A,
 \qquad
 A\left[1-\left(\frac{65}{289}\right)^m\right]
 <\operatorname{Re}\Phi_h(r_j)
 <A\left[1+\left(\frac{65}{289}\right)^m\right]
 \quad(1\leq j\leq4).
 \tag{QCV.19}
\]
This estimate has not assumed that an arithmetic zero has height at
least two. That fact is proved from the original theta kernel in
(LHT.1)--(LHT.18) below. Those equations give
\(\operatorname{Re}g(s)>97/2052\) for
\(0\leq\operatorname{Re}s\leq1\), \(|\operatorname{Im}s|\leq2\).
Consequently every actual zero (QCV.1) of the original \(g=2\xi\)
has \(\gamma>2\), and (QCV.19) applies to every actual offcritical
quartet with its full order. This use of the low-height calculation
does not assert that offcritical quartets exist.

\subsection{Every ordered tensor critical point, collision and full local rank}

For \(k\geq1\), keep the ordered tensor phase on the original
coordinates
\[
 \Phi_{h,k}(s_1,\ldots,s_k)=\sum_{i=1}^k\Phi_h(s_i),
 \quad E_{h,k}=\mathbb C[s_1,\ldots,s_k]/(h(s_1),\ldots,h(s_k)).
 \tag{QCV.20}
\]
Its critical points are precisely the ordered tuples
\(\mathbf r=(r_{j_1},\ldots,r_{j_k})\), \(j_i\in\{1,2,3,4\}\),
since its \(i\)-th derivative is \(h(s_i)\). Write \(n_j\) for
the number of occurrences of \(r_j\); thus \(n_j\geq0\) and
\(\sum_jn_j=k\). Evaluation gives exactly
\[
 \Phi_{h,k}(\mathbf r)=kA+(n_1+n_2-n_3-n_4)U
                           +i(n_1-n_2-n_3+n_4)V.
 \tag{QCV.21}
\]
With \(a=n_1+n_2\) and \(b=n_1+n_4\), these values are
\[
 c_{k,a,b}=kA+(2a-k)U+i(2b-k)V,
                  \qquad 0\leq a,b\leq k.
 \tag{QCV.22}
\]
Every pair \((a,b)\) occurs: choose a set of \(a\) positions for a
positive \(U\)-sign and, independently, a set of \(b\) positions
for a positive \(V\)-sign. At each position the pair of signs picks
exactly one of \(r_1,r_2,r_3,r_4\) in (QCV.8). This also proves
that the number of ordered tuples giving \(c_{k,a,b}\) is exactly
\[
 N_{k,a,b}=\binom ka\binom kb
 =\sum_{x=\max(0,a+b-k)}^{\min(a,b)}
   \frac{k!}{x!(a-x)!(k-a-b+x)!(b-x)!}.
 \tag{QCV.23}
\]
The sum follows by setting
\((n_1,n_2,n_3,n_4)=(x,a-x,k-a-b+x,b-x)\) and counting
permutations. Since \(U,V>0\), two values in (QCV.22) are equal
if and only if both their indices \(a,b\) agree. Thus there are
exactly \((k+1)^2\) distinct critical values. All collisions among
different ordered tuples are the ones counted in (QCV.23).

At an individual tuple set \(y_i=s_i-r_{j_i}\). Polynomial Chinese
remaindering gives the exact decomposition
\[
 E_{h,k}\simeq
 \bigoplus_{(j_1,\ldots,j_k)\in\{1,2,3,4\}^k}
 \mathbb C[y_1,\ldots,y_k]/(y_1^m,\ldots,y_k^m).
 \tag{QCV.24}
\]
For completeness, the ideals \(((s-r_j)^m)\) are pairwise
comaximal: distinct \(r_j\) give coprime polynomials, and the
Euclidean algorithm supplies polynomials expressing \(1\) as a
sum of multiples of each pair. The Chinese-remainder map is
evaluation of every full Taylor jet \(s=r_j+y\) modulo \(y^m\);
its kernel is the product ideal \((h)\), and the same Bezout
identities construct the inverse idempotents. Tensoring these
finite-dimensional isomorphisms in the unchanged order gives
(QCV.24). The monomials \(\prod_i y_i^{e_i}\), \(0\leq e_i<m\),
form a basis of each summand, so its dimension is \(m^k\).

The summand belonging to a coincident critical value \(c_{k,a,b}\)
therefore has dimension
\[
 d_{k,a,b}=m^k\binom ka\binom kb,
 \qquad \sum_{a,b=0}^kd_{k,a,b}=(4m)^k.
 \tag{QCV.25}
\]
Grouping by equal value is a direct sum of the original full local
algebras, and does not identify or quotient their jets. In each
individual summand the original arithmetic tensor sum is still
\[
 A_{h,k}=\left(\sum_{i=1}^kr_{j_i}\right)I+\sum_{i=1}^kM_{y_i},
 \quad
 \exp(tA_{h,k})=e^{t\sum_i r_{j_i}}
       \prod_{i=1}^k\sum_{e=0}^{m-1}\frac{t^eM_{y_i}^e}{e!}.
 \tag{QCV.26}
\]
All nilpotent powers remain, and the scalar \(\sum_i r_{j_i}\)
in this arithmetic action is distinct data from the phase value
\(\sum_i\Phi_h(r_{j_i})\). Both are calculated on the same
unchanged tuple; no equality between them is asserted.

For an actual zero of \(g\), (QCV.19) and (QCV.22) finally give
\[
 \boxed{\begin{gathered}
 \operatorname{Re}c_{k,a,b}
       \geq k(A-U)
       >kA\left[1-\left(\frac{65}{289}\right)^m\right]>0,\\
 |c_{k,a,b}|\geq\operatorname{Re}c_{k,a,b}
       >\frac{k}{2}(\delta^2+\gamma^2)^{2m}
                       \left[1-\left(\frac{65}{289}\right)^m\right].
 \end{gathered}}
 \tag{QCV.27}
\]
Every inequality uses the original constant \(A\), with the explicit
lower bound for \(A\) inserted only in the second line. In particular
no original ordered tensor critical value is zero. This proves the
nonvanishing of the phase constants required at those precise
vertices of the programme; it does not identify the arithmetic
tensor action with a Frobenius action, and it proves neither the
existence nor the absence of an offcritical zero by itself.
